% GENERATED FILE. Edit canonical lessons, then run npm run generate:books.
% canonical-source-hash: fnv1a64:3a644088456519c4
% canonical-lessons: MR-C281-jayla-pahije, MR-C281-karan-mule, MR-C281-jnan, MR-W281-letter-nya, MR-R281-advice-and-reasons

\chapter{Must, Must Not, and Because}
\label{ch:mr-must-must-not-and-because}
\hlchaptermodality{281}

\begin{chapteropening}
\textbf{By the end of this chapter:} \emph{I can say what someone must do (jāylā pāhije) and shouldn't do (jāū nayes), give a reason with -muḷe and kāraṇ, and read and write \mr{ञ} in jñān.}

Complete MR-R281-advice-and-reasons: give advice and reasons, after meeting \mr{ज्ञान} and its letter \mr{ञ}.
\end{chapteropening}

\section[\texorpdfstring{\mr{मला} \mr{जायला} \mr{पाहिजे} (malā jāylā pāhije)}{मला जायला पाहिजे (malā jāylā pāhije)}]{\mr{मला} \mr{जायला} \mr{पाहिजे} (malā jāylā pāhije) — must, and must not}

\label{lesson:MR-C281-jayla-pahije}



\begin{quote}
Say \textquotedblleft{}to go for a walk\textquotedblright{}: \emph{phirāylā jāṇe}. \textbf{\mr{फिरायला}} is \textquotedblleft{}for walking\textquotedblright{}, and it is the form this lesson needs.
\end{quote}

\subsection*{You'll want to know: -\mr{ायला} \mr{पाहिजे}, 'must, should'}
\begin{quote}
\textbf{\mr{मला} \mr{आता} \mr{घरी} \mr{जायला} \mr{पाहिजे}.}
\end{quote}

\begin{quote}
\emph{malā ātā gharī jāylā pāhije.} — \textquotedblleft{}I must go home now.\textquotedblright{}
\end{quote}

Take \textbf{-\mr{णे}} off and put \textbf{-\mr{ायला}} on: \textbf{\mr{जाणे}} $\to$ \textbf{\mr{जायला}} (\emph{jāylā}), \textbf{\mr{करणे}} $\to$ \textbf{\mr{करायला}} (\emph{karāylā}), \textbf{\mr{येणे}} $\to$ \textbf{\mr{यायला}} (\emph{yāylā}), \textbf{\mr{वाचणे}} $\to$ \textbf{\mr{वाचायला}}. Then \textbf{\mr{पाहिजे}} (\emph{pāhije}), which does not change. The one who must stands in the dative, as with \textbf{\mr{आवडणे}}: \textbf{\mr{मला}}, \textbf{\mr{तुला}} (\emph{tulā}, \textquotedblleft{}to you, a friend\textquotedblright{}), \textbf{\mr{त्याला}} (\emph{tyālā}, \textquotedblleft{}to him\textquotedblright{}), \textbf{\mr{तिला}} (\emph{tilā}, \textquotedblleft{}to her\textquotedblright{}).

\textbf{\mr{हवं}}, the \textbf{\mr{हवा}} you know, does the same work: \textbf{\mr{मला} \mr{जायला} \mr{हवं}} (\emph{malā jāylā havaṁ}).

\subsection*{You'll want to know: -\mr{ऊ} \mr{नये}, 'should not'}
\begin{quote}
\textbf{\mr{तू} \mr{रात्री} \mr{जाऊ} \mr{नयेस}.}
\end{quote}

\begin{quote}
\emph{tū rātrī jāū nayes.} — \textquotedblleft{}You shouldn't go at night.\textquotedblright{}
\end{quote}

The \textbf{-\mr{ऊ}} of \textbf{\mr{आम्ही} \mr{जाऊ}} with \textbf{\mr{नये}}, \textquotedblleft{}should not\textquotedblright{}, which takes \textbf{-\mr{स}} after \textbf{\mr{तू}}. Said to no one in particular: \textbf{\mr{असं} \mr{करू} \mr{नये}} (\emph{asaṁ karū naye}), \textquotedblleft{}one shouldn't do that\textquotedblright{}. A flat \textquotedblleft{}don't\textquotedblright{} uses the \textbf{\mr{नको}} you know: \textbf{\mr{जाऊ} \mr{नकोस}} (\emph{jāū nakos}), \textquotedblleft{}don't go\textquotedblright{}.

\subsection*{Your turn}
Say these aloud:

\begin{itemize}
\raggedright
  \item \textquotedblleft{}I must go home now\textquotedblright{} — \emph{malā ātā gharī jāylā pāhije}
  \item \textquotedblleft{}He should read this book\textquotedblright{} — \emph{tyālā he pustak vācāylā pāhije}
  \item \textquotedblleft{}You shouldn't go at night\textquotedblright{} — \emph{tū rātrī jāū nayes}
\end{itemize}

\begin{tcolorbox}[breakable,colback=teal!4,colframe=teal!35!black,title={Before you move on}]
\textquotedblleft{}I must go\textquotedblright{}? (\textbf{\mr{मला} \mr{जायला} \mr{पाहिजे}}.) Does \textbf{\mr{पाहिजे}} change for \textquotedblleft{}she\textquotedblright{}? (No: \textbf{\mr{तिला} \mr{जायला} \mr{पाहिजे}}.) \textquotedblleft{}You shouldn't go\textquotedblright{}? (\textbf{\mr{तू} \mr{जाऊ} \mr{नयेस}}.)
\end{tcolorbox}
\section[\texorpdfstring{-\mr{मुळे} (-muḷe)}{-मुळे (-muḷe)}]{-\mr{मुळे} (-muḷe) — because of, and the reason in full}

\label{lesson:MR-C281-karan-mule}



\begin{quote}
Say \textquotedblleft{}because of that\textquotedblright{}. (\emph{tyāmuḷe}.) Its second half is the ending this lesson takes apart.
\end{quote}

\subsection*{You'll want to know: -\mr{मुळे}, 'because of'}
\begin{quote}
\textbf{\mr{पावसामुळे} \mr{गाडी} \mr{उशिरा} \mr{आली}.}
\end{quote}

\begin{quote}
\emph{pāvsāmuḷe gāḍī uśirā ālī.} — \textquotedblleft{}Because of the rain, the train came late.\textquotedblright{}
\end{quote}

\textbf{-\mr{मुळे}} goes on a noun in its bent form. \textbf{\mr{पाऊस}} bends to \textbf{\mr{पावसा}-}, the same as in \textbf{\mr{पावसाळा}}: \textbf{\mr{पावसामुळे}} (\emph{pāvsāmuḷe}). \textbf{\mr{उशिरा}} (\emph{uśirā}, \textquotedblleft{}late\textquotedblright{}) is \textbf{\mr{उशीर}} made into \textquotedblleft{}how\textquotedblright{}. On a pronoun: \textbf{\mr{तुमच्यामुळे}} (\emph{tumcyāmuḷe}), \textquotedblleft{}because of you\textquotedblright{}; and \textbf{\mr{त्यामुळे}}, \textquotedblleft{}because of that\textquotedblright{}.

\subsection*{You'll want to know: the reason in full}
\begin{quote}
\textbf{\mr{मी} \mr{काल} \mr{आलो} \mr{नाही}, \mr{कारण} \mr{मी} \mr{आजारी} \mr{होतो}.}
\end{quote}

\begin{quote}
\emph{mī kāl ālo nāhī, kāraṇ mī ājārī hoto.} — \textquotedblleft{}I didn't come yesterday, because I was ill.\textquotedblright{}
\end{quote}

When someone asks \textbf{\mr{तुम्ही} \mr{का} \mr{आला} \mr{नाहीत}?} (\emph{tumhī kā ālā nāhīt?}), \textquotedblleft{}Why didn't you come?\textquotedblright{}, answer with the event and then \textbf{\mr{कारण}} and the reason, a whole sentence in the past. Or the reason first and then \textbf{\mr{म्हणून}}: \textbf{\mr{मी} \mr{आजारी} \mr{होतो}, \mr{म्हणून} \mr{मी} \mr{आलो} \mr{नाही}} (\emph{mī ājārī hoto, mhaṇūn mī ālo nāhī}).

\subsection*{Your turn}
Say these aloud:

\begin{itemize}
\raggedright
  \item \textquotedblleft{}Because of the rain, the train came late\textquotedblright{} — \emph{pāvsāmuḷe gāḍī uśirā ālī}
  \item \textquotedblleft{}Why didn't you come?\textquotedblright{} — \emph{tumhī kā ālā nāhīt?}
  \item \textquotedblleft{}I didn't come, because I was ill\textquotedblright{} — \emph{mī ālo nāhī, kāraṇ mī ājārī hoto}
\end{itemize}

\begin{tcolorbox}[breakable,colback=teal!4,colframe=teal!35!black,title={Before you move on}]
\textquotedblleft{}Because of the rain\textquotedblright{}? (\textbf{\mr{पावसामुळे}}.) \textquotedblleft{}Because of you\textquotedblright{}? (\textbf{\mr{तुमच्यामुळे}}.) Which takes a whole sentence, \textbf{-\mr{मुळे}} or \textbf{\mr{कारण}}? (\textbf{\mr{कारण}}.)
\end{tcolorbox}
\section[\texorpdfstring{\mr{ज्ञान} (jñān)}{ज्ञान (jñān)}]{\mr{ज्ञान} (jñān) — knowledge}

\label{lesson:MR-C281-jnan}



\begin{quote}
Before the new one: say \textquotedblleft{}a book\textquotedblright{} and \textquotedblleft{}to grow\textquotedblright{}. (\emph{pustak}, \emph{vāḍhṇe}.)
\end{quote}

\subsection*{You'll want to know: \mr{ज्ञान}}
\textbf{\mr{ज्ञान}} — \emph{jñān} — \textquotedblleft{}knowledge\textquotedblright{}.

\begin{quote}
\textbf{\mr{पुस्तकांमुळे} \mr{ज्ञान} \mr{वाढतं}.}
\end{quote}

\begin{quote}
\emph{pustakāmmuḷe jñān vāḍhtaṁ.} — \textquotedblleft{}Knowledge grows because of books.\textquotedblright{}
\end{quote}

The first block, \textbf{\mr{ज्ञ}}, is said \emph{dny} in Marathi: \textbf{\mr{ज्ञान}} is said \emph{dnyān}. The romanization keeps the letters, \emph{jñ}. \textbf{\mr{ज्ञान} \mr{महत्त्वाचं} \mr{आहे}} (\emph{jñān mahattvācaṁ āhe}): knowledge is important.

\begin{cousinweb}
\textbf{\mr{ज्ञान}} (\emph{jñān}) is a tatsama: Sanskrit \emph{jñāna}, \textquotedblleft{}knowledge\textquotedblright{}, on the root \emph{jñā}, \textquotedblleft{}to know\textquotedblright{} (CDIAL 5281). Under the same entry CDIAL lists an inherited Marathi form, \emph{jāṇ}, in the narrower sense \textquotedblleft{}acknowledgement of favours\textquotedblright{}; \textbf{\mr{ज्ञान}} is the learned word, taken whole from Sanskrit.
\end{cousinweb}

\subsection*{Your turn}
Say these aloud:

\begin{itemize}
\raggedright
  \item \emph{jñān}, said \emph{dnyān}
  \item \textquotedblleft{}Knowledge grows because of books\textquotedblright{} — \emph{pustakāmmuḷe jñān vāḍhtaṁ}
  \item \textquotedblleft{}Knowledge is important\textquotedblright{} — \emph{jñān mahattvācaṁ āhe}
\end{itemize}

\begin{tcolorbox}[breakable,colback=teal!4,colframe=teal!35!black,title={Before you move on}]
What does \textbf{\mr{ज्ञान}} mean? (\textquotedblleft{}Knowledge\textquotedblright{}.) How is \textbf{\mr{ज्ञ}} said? (\emph{dny}.)
\end{tcolorbox}
\section[\texorpdfstring{\mr{ञ} (ña)}{ञ (ña)}]{\mr{ञ} (ña) — the letter hidden inside \mr{ज्ञ}}

\label{lesson:MR-W281-letter-nya}



\begin{quote}
Say \textbf{\mr{ज्ञान}}, \textquotedblleft{}knowledge\textquotedblright{}. You have read its first block, \textbf{\mr{ज्ञ}}, as one shape. Name the letter it starts with: \textbf{\mr{ज}}. Then the mark that silences a letter's vowel: \textbf{\mr{्}}.
\end{quote}

\subsection*{Script}
\begin{quote}
\mr{ञ}
\end{quote}

\textbf{\mr{ञ}} is \emph{ña}, a nasal made where \emph{j} is made. \textbf{\mr{ज्ञ}} is \textbf{\mr{ज}}, the \textbf{\mr{्}}, and \textbf{\mr{ञ}}, joined into one shape; Marathi says the pair \emph{dny}. On its own \textbf{\mr{ञ}} is rare in Marathi words; you will meet it inside \textbf{\mr{ज्ञ}}, in \textbf{\mr{ज्ञान}}.

\subsection*{Writing: \mr{ञ} — guided copy}
\hlblockfigure{\detokenize{figures/MR-W281-letter-nya-filmstrip.pdf}}{How \mr{ञ} is written, stroke by stroke}

Keep \textbf{\mr{ञ}} in front of you for this first pass.

\begin{itemize}
\raggedright
  \item \textbf{1.} start at the open upper-left tip and curve clockwise around the bowl to the open lower-left tip
  \item \textbf{2.} lift at the bowl's right junction, sweep right through the shoulder, then curve upward along the right side to the headline
  \item \textbf{3.} lift at the lower right-side junction and descend the short stem below the bowl
  \item \textbf{4.} lift and draw the top shirorekhā left-to-right
\end{itemize}

\textbf{Pen lifts: 3.} This is one attested teaching order; hands and teachers vary. Source: Opiaterein, ‘Deva-\mr{ञ}-order.gif’, strokes 1–4, Wikimedia Commons, 10 May 2009.

Copy \textbf{\mr{ञ}} twice beside the model, saying \emph{ña} as you finish each one. Then write \textbf{\mr{ज्ञान}} once and point at the \textbf{\mr{ञ}} inside its first block.

\begin{tcolorbox}[breakable,colback=teal!4,colframe=teal!35!black,title={Before you move on}]
Which letter hides inside \textbf{\mr{ज्ञ}}? (\textbf{\mr{ञ}}, \emph{ña}.) Which word holds it? (\textbf{\mr{ज्ञान}}, \emph{jñān}, said \emph{dnyān}.)
\end{tcolorbox}
\section[(review)]{Review — advice and reasons}

\label{lesson:MR-R281-advice-and-reasons}



\begin{quote}
Say \textquotedblleft{}I think it will rain tomorrow, so I will stay at home\textquotedblright{}. (\emph{malā vāṭtaṁ kī udyā pāūs paḍel, mhaṇūn mī gharī rāhīn}.) Then say \textbf{\mr{ऋतू}}, and \textbf{\mr{ज्ञान}} with its \emph{dny}.
\end{quote}

\subsection*{Your turn}
Say these aloud:

\begin{itemize}
\raggedright
  \item I must go home now, because it is late.
  \item You shouldn't go at night.
  \item Because of the rain, the train came late.
  \item Why didn't you come? I didn't come, because I was ill.
  \item In my opinion knowledge grows because of books.
\end{itemize}

\begin{tcolorbox}[breakable,colback=teal!4,colframe=teal!35!black,title={Before you move on}]
\textquotedblleft{}Must\textquotedblright{}? (\textbf{-\mr{ायला} \mr{पाहिजे}}.) \textquotedblleft{}Should not\textquotedblright{}? (\textbf{-\mr{ऊ} \mr{नये}}.) \textquotedblleft{}Because of\textquotedblright{}? (\textbf{-\mr{मुळे}}.) Which letter hides in \textbf{\mr{ज्ञ}}? (\textbf{\mr{ञ}}.) Which starts \textbf{\mr{ऋतू}}? (\textbf{\mr{ऋ}}.)
\end{tcolorbox}
